% !TEX root = ../../tree_notes.tex
% =====================================================================
%  1.2  Depth, Height and Size                           sec:c01-anatomy
%  Exports: def:c01-depth  def:c01-height  def:c01-size
%           thm:c01-height-recursion  lst:c01-height  lst:c01-measure
%           prob:c01-max-depth  prob:c01-min-depth  prob:c01-count-nodes-avg
% =====================================================================
\section{Depth, Height and Size}\label{sec:c01-anatomy}

\begin{previouslybox}
In \secref{c01-what-is-tree} we earned the family vocabulary (root, parent, child,
leaf, subtree) and the two theorems every tree algorithm leans on: $n$ nodes $\to$
$n-1$ edges, and one unique path per pair. We also met the parent-guarded DFS of
\coderef{c01-build}. Keep both in hand -- this section attaches \emph{numbers} to
that vocabulary.

\begin{enumerate}[nosep, leftmargin=1.6em]
  \item A tree on $n$ nodes has \underline{\hspace{0.8cm}} edges, and \emph{any} two
        nodes are joined by \underline{\hspace{0.9cm}} simple path(s).
        \emph{(Ans: $n-1$; exactly one -- \thmref{c01-unique-path}.)}
  \item In the running example of \figref{c01-anatomy}, is $D$ an ancestor of $B$, a
        descendant of $B$, or a sibling of $C$?
        \emph{(Ans: a descendant of $B$; $D$ and $C$ are not siblings -- their parents
        differ.)}
  \item Why does the DFS of \coderef{c01-build} pass the parent $p$ to every recursive
        call? \emph{(Ans: adjacency edges are undirected, so without skipping $p$ the
        walk bounces back and never terminates.)}
\end{enumerate}
\end{previouslybox}

\begin{topicintro}
Three numbers describe every node of a rooted tree: how far it is from the root
(\term{depth}), how far the deepest leaf below it is (\term{height}), and how many
nodes hang beneath it (\term{size}). Nearly every tree algorithm is "compute these
bottom-up, then combine them," so they are the first tools to make automatic.

\textbf{Prerequisites.} \secref{c01-what-is-tree} (parent, child, subtree) and the DFS of
\coderef{c01-build}.

\textbf{After this section you can}
\begin{itemize}[nosep]
  \item state $\depth$, $\height$ and $\size$ for every node of a small tree by hand;
  \item compute all three for the whole tree in a single $O(n)$ DFS;
  \item write the height recursion correctly under both the edges and nodes conventions;
  \item avoid the stack overflow that a skewed tree causes on recursive solutions.
\end{itemize}
\end{topicintro}

\begin{recognitionbox}[Pattern card -- top-down vs bottom-up measurement]
\reckey{``depth / level / distance from root'' (carry a parameter \emph{down});
``height / deepest / how many below / subtree total'' (return a value \emph{up}).
Almost every ``compute something for every node'' problem is one of these two.}
\recuse{
  \item the answer at a node depends only on its \emph{ancestors} $\longrightarrow$ top-down
        (depth, path prefixes, root-to-leaf sums);
  \item the answer at a node depends on its \emph{descendants} $\longrightarrow$ bottom-up
        (height, size, subtree sums, any $\max/\min/\mathrm{sum}$ over a subtree);
  \item both at once -- one DFS can carry depth down and return height/size up
        (\coderef{c01-measure}).
}
\recnotuse{
  \item the quantity needs \emph{sibling} subtrees combined at a node and a global best --
        that is diameter-style logic (\chref{c09}) or tree DP (\chref{c10});
  \item queries mix updates and subtree reads -- Euler tour + a data structure (\chref{c04},
        \chref{c13});
  \item a plain traversal suffices (visit/print/count all) -- BFS/DFS of \chref{c03}.
}
One-line test: \emph{``does the node's value flow from the root, or bubble up from the
leaves?'' The direction picks the recursion.}
\end{recognitionbox}

\subsection{Intuition}

Root the tree of \figref{c01-anatomy} at $A$ and imagine two rulers. One hangs from the
root downward and reads the \emph{depth} of each node; the other stands on each node and
reaches down to the deepest leaf below it, reading the \emph{height}.

\begin{figure}[H]
\centering
\begin{tikzpicture}[diagram]
  \node[troot] (a) at (0,0) {A};
  \node[tnode] (b) at (-2,-1.5) {B};
  \node[tnode] (c) at (2,-1.5) {C};
  \node[tnode] (d) at (-2.9,-3) {D};
  \node[tleaf] (e) at (-1.1,-3) {E};
  \node[tleaf] (f) at (2,-3) {F};
  \draw[tedge] (a) -- (b); \draw[tedge] (a) -- (c);
  \draw[tedge] (b) -- (d); \draw[tedge] (b) -- (e); \draw[tedge] (c) -- (f);
  % depth labels on the left
  \node[annot, text=primaryblue] at (-3.7,-1.5) {$\depth{=}1$};
  \node[annot, text=primaryblue] at (-3.9,-3) {$\depth{=}2$};
  % height labels on the right
  \node[annot, text=successgreen] at (3.6,-3) {$h{=}0$};
  \node[annot, text=successgreen] at (3.3,-1.5) {$h{=}1$};
  \node[annot, text=accentorange] at (3.3,0.1) {$\height(T){=}2$};
\end{tikzpicture}
\caption{Blue shows \emph{depth} (measured down from the root, in edges); green shows
\emph{height} of a few nodes (edges to the deepest leaf below them). The whole tree has
height $2$, the height of its root $A$. Leaves have height $0$.}\label{fig:c01-depth-height}
\end{figure}

\vsbox{Depth vs Height}{Depth $\depth(v)$}{Height $\height(v)$}{edges from the
\emph{root} down to $v$; shared by the walk from the top}{edges from $v$ down to its
deepest \emph{leaf}; computed from the bottom up}

\subsection{Definitions and Notation}

\begin{definition}{Depth and level}{c01-depth}
The \term{depth} $\depth(v)$ of a node is the number of edges on the path from the root
to $v$; the root has $\depth=0$. The \term{level} of a node is its depth (some books call
the root level $1$; we always start levels at $0$). Siblings share a depth.
\end{definition}

\noindent\textbf{In plain words.} Depth answers ``how many steps down from the boss am
I?'' -- the boss herself counts as $0$. Think generations in a family photo: same
generation, same depth, always.

\begin{definition}{Height}{c01-height}
The \term{height} $\height(v)$ of a node is the number of edges on the longest path from
$v$ down to a leaf in its subtree; a leaf has $\height=0$ and the empty child has
$\height(\nil)=-1$. The height of the tree is $\height(T)=\height(\text{root})$.
\end{definition}

\noindent\textbf{In plain words.} Height answers the opposite question: ``looking
\emph{down}, how long is the deepest chain I manage?'' A leaf manages nothing, so $0$.
The odd-looking $\height(\nil)=-1$ is not philosophy -- it is the one convention that
makes the $\max$ formula of \thmref{c01-height-recursion} compute leaves correctly all by
itself, with no special case in your code.

\begin{definition}{Size}{c01-size}
The \term{size} $\size(v)$ is the number of nodes in $\subtree(v)$, counting $v$
itself; so $\size(\nil)=0$ and $\sum$ of nothing is $0$.
\end{definition}

\noindent\textbf{In plain words.} Size is a headcount inside your own department -- and
you count yourself. An empty department has nobody to count: $\size(\nil)=0$.

\subsection{The Bottom-Up Recurrences}

Height and size both "bubble up" from the leaves: a node reads the answers its children
computed and folds them together. This is the pattern behind half the book.

\begin{theorem}{Height and Size Recursion}{c01-height-recursion}
With $\height(\nil)=-1$ and $\size(\nil)=0$, for any node $v$
\[
  \height(v) = 1 + \max_{c \in \children(v)} \height(c),
  \qquad
  \size(v) = 1 + \sum_{c \in \children(v)} \size(c).
\]
A leaf ($\children(v)=\emptyset$) gets $\height=0$ and $\size=1$.
\end{theorem}
\begin{proofbox}
The longest downward path from $v$ must leave $v$ through exactly one child, then follow
the longest path from that child; that is $1+\max_c \height(c)$, and if there is no child
the max over the empty set is $\height(\nil)=-1$, giving $\height(v)=0$. Every node of
$\subtree(v)$ is either $v$ itself or a node of exactly one child's subtree, so the counts
add to $1+\sum_c \size(c)$. Both are well-founded because a tree has no cycles, so
recursion reaches a leaf in finitely many steps.
\end{proofbox}
\begin{insightbox}
Depth is a \emph{top-down} quantity (carry it into the recursion as a parameter), while
height and size are \emph{bottom-up} (returned from it). One DFS can do all three: push
$\depth$ in, return $\height$ and $\size$ out.
\end{insightbox}

\usedin{\chref{c06} (AVL balance factors), \chref{c09} (diameter) and \chref{c18} (centroid
decomposition)}{this same "compute a value per node bottom-up, combine children" move
returns again and again with a different local formula.}

\subsection{One DFS that Measures Everything}

\begin{lstlisting}[caption={Height of a binary tree (height in edges, height(null) = -1)}, label={lst:c01-height}]
// height in EDGES: null -> -1, leaf -> 0.  (See the note below for the nodes version.)
int height(TreeNode* r) {
    if (r == nullptr) return -1;
    return 1 + max(height(r->left), height(r->right));
}
\end{lstlisting}

\begin{lstlisting}[caption={One DFS that fills depth, height and size for a general tree}, label={lst:c01-measure}]
vector<vector<int>> adj;
vector<int> depth_, height_, size_;

void measure(int v, int p, int d) {         // p = parent, d = depth of v
    depth_[v]  = d;
    size_[v]   = 1;                          // count v itself
    height_[v] = -1;                         // a leaf ends at -1 + 1 = 0
    for (int to : adj[v]) {
        if (to == p) continue;               // parent guard (see 1.1)
        measure(to, v, d + 1);               // push depth down
        size_[v]   += size_[to];             // pull size up
        height_[v]  = max(height_[v], height_[to]);
    }
    height_[v] += 1;                         // add the edge to the deepest child
}
\end{lstlisting}

\begin{complexitybox}
Each node is visited once and each edge twice, so the DFS is $\bigO{n}$ time and
$\bigO{n}$ memory (arrays plus a recursion stack up to the height of the tree). On a
path of $n$ nodes the stack is $\bigO{n}$ deep -- see the pitfall below.
\end{complexitybox}

\begin{notebox}[Edges vs nodes: an off-by-one you must choose]
This book counts height/depth in \textbf{edges}: $\height(\nil)=-1$, a lone node has
height $0$. LeetCode's \emph{Maximum Depth} (\coderef{c01-max-depth-soln} below) counts in
\textbf{nodes}: $\height(\nil)=0$, a lone node has depth $1$. The two answers differ by
exactly one. Read the statement, pick the base case, and never mix them in one program.
\end{notebox}

\subsection{Worked Examples}

\begin{example}{Reading all three numbers}{c01-anatomy-ex}
In \figref{c01-depth-height}: $D$ has $\depth 2, \height 0, \size 1$; $B$ has
$\depth 1, \height 1, \size 3$; $A$ has $\depth 0, \height 2, \size 6$. Note
$\size(A)=6$ equals the node count, and $\height(A)$ is the tallest root-to-leaf run of
edges ($A\to B\to D$ or $A\to B\to E$).
\end{example}

\begin{tracebox}
Running \lstname\ of \coderef{c01-measure} on a path $0-1-2$ rooted at $0$ fills
depth $\{0,1,2\}$, height $\{2,1,0\}$ and size $\{3,2,1\}$. Height grows from the far
leaf back to the root, size collapses onto the root -- the two recurrences running in
opposite directions along the same walk.
\end{tracebox}

\subsection{Variations and Twists}

\begin{itemize}
  \item \textbf{Direction flip.} Every measurement here has a top-down twin: depth is
        carried \emph{in}, height is returned \emph{out}. LC 1123 (LCA of deepest
        leaves) is the bottom-up height recursion with an LCA twist -- \chref{c11}.
        \hfill\lv{3}
  \item \textbf{Count the level, not the edge.} Replace $\max$ by "which level has the
        most nodes" and the same DFS becomes level statistics -- width
        questions (\textit{Maximum Width of Binary Tree}, LC 662) index nodes by
        position and run on the same rails. \hfill\lv{3}
  \item \textbf{Two quantities, one DFS.} The stretch problem below (LC 2265) needs
        \emph{size and sum} returned together -- the exact shape of
        \coderef{c01-measure}, one field richer. This is the whole gateway into
        tree DP (\chref{c10}): a state per node, folded bottom-up. \hfill\lv{4}
  \item \textbf{The iterative demand.} On judges with $n \le 10^{5}$ chain-shaped
        trees every entry here needs a stack machine version; see the hardening note
        under the stretch solution. \hfill\lv{3}--\lv{4}
  \item \textbf{Disguise.} "Time needed to inform all employees" (LC 1376) is a
        general tree with depths; "deepest leaves sum" is depth-bucketed accumulation.
        Spotting \emph{it is just depth} is the entire solve.
\end{itemize}

\subsection{Common Pitfalls}

\begin{warningbox}[Recursion on a skewed tree]
A tree that is really a path of $10^5$ nodes makes the DFS recurse $10^5$ deep and can
overflow the call stack. Prefer an explicit stack (iterative post-order) or raise the
stack limit when the judge allows. Every recursive tree routine in this book gets an
iterative twin for exactly this reason.
\end{warningbox}

\begin{warningbox}[Off-by-one on the base case]
Writing $\height(\nil)=0$ when the problem wants edges (or vice-versa) shifts every
answer by one. Decide the convention \emph{first}, set the null base case to match, and
test on a single-node tree -- the fastest way to catch it.
\end{warningbox}

\smallskip\noindent\textbf{The traps as code.}

\begin{lstlisting}[caption={\textcolor{warningred}{WRONG -- mixed conventions:} LC-104 statement (nodes) with the edges base case}, basicstyle=\ttfamily\small]
int maxDepth(TreeNode* r) {
  if (!r) return -1;   // edges base ...
  return 1 + max(maxDepth(r->left),
                 maxDepth(r->right)); // ... +1 per node!
}
\end{lstlisting}
\begin{lstlisting}[caption={\textcolor{successgreen}{RIGHT -- pick one convention and keep it (nodes here)}}, basicstyle=\ttfamily\small, breaklines=true]
int maxDepth(TreeNode* r) {
  if (!r) return 0;                      // nodes: null = 0
  return 1 + max(maxDepth(r->left), maxDepth(r->right));
}
\end{lstlisting}

\subsection{Problem Lab (fully solved)}

\begin{attemptgate}
\textbf{Maximum Depth} is $\lv1$ -- write it cold in under 3 minutes, both
conventions from memory. \textbf{Minimum Depth} is the one that eats
interviews: get a sheet of paper, draw the three trees in its
\hyperref[fig:c01-min-depth-trap]{trap figure}, and predict the answer
\emph{before} reading a line of code.
\end{attemptgate}

\solvedproblem{c01-max-depth}{Maximum Depth of Binary Tree}{LeetCode 104}
  {https://leetcode.com/problems/maximum-depth-of-binary-tree/}{LC Easy L1}
  {recursion, height of a binary tree}

\begin{problemstatement}[Maximum Depth]
Given the root of a binary tree, return the number of \emph{nodes} on the longest path
from the root down to the farthest leaf. A single node has depth $1$; an empty tree has
depth $0$.
\end{problemstatement}

\solstep{Understanding the Problem}
This is the height of the tree counted in \emph{nodes}, not edges -- LeetCode's
convention. So the answer is the edge-height of \thmref{c01-height-recursion} plus one.
The natural formulation is recursive: the deepest path through a node is one more than
the deeper of its two children.

\solstep{Key Observations}
\begin{observation}
Counting depth in nodes: $\depth(v) = 1 + \max(\depth(\text{left}),\depth(\text{right}))$,
with the base case $\depth(\nil)=0$.
\end{observation}
\begin{observation}
This mirrors the height recursion exactly; only the base case changes ($0$, not $-1$).
\end{observation}

\solstep{Solution}
Recurse, returning $0$ for a null child and $1+\max(\text{left},\text{right})$ otherwise.
\begin{figure}[H]
\centering
\begin{forest} snaptree [3 [9 [4] [1]] [,phantom]] \end{forest}
\caption{The longest root-to-leaf path is $3\to 9\to 4$ (or $3\to 9\to 1$), three
\emph{nodes}, so the answer is $3$. Counting edges instead would give $2$.}\label{fig:c01-max-depth}
\end{figure}

\solstep{Proof of Correctness}
\begin{proofbox}[Induction on $\size(v)$]
Base: $v=\nil$, the code returns $0$, correct. Step: assume the calls on the children
return their correct node-depths. The deepest path from $v$ goes to the deeper child and
then adds the single node $v$, so $1+\max$ is exactly the longest root-to-leaf node count
at $v$ -- the same recurrence the theorem states, shifted by the base case.
\end{proofbox}

\solstep{Implementation}
\begin{lstlisting}[caption={Maximum Depth (LeetCode 104)}, label={lst:c01-max-depth-soln}]
#include <bits/stdc++.h>
using namespace std;

struct TreeNode {
    int val;
    TreeNode *left, *right;
    TreeNode(int x = 0) : val(x), left(nullptr), right(nullptr) {}
};

class Solution {
public:
    int maxDepth(TreeNode* root) {
        if (root == nullptr) return 0;                       // nodes convention
        return 1 + max(maxDepth(root->left), maxDepth(root->right));
    }
};
\end{lstlisting}

\solstep{Complexity}
\begin{complexitybox}
Time $\bigO{n}$ (every node visited once). Memory $\bigO{h}$ recursion stack, worst-case
$\bigO{n}$ on a skewed tree; a level-order (BFS) count of how many levels there are gives
the same answer iteratively with $\bigO{w}$ queue space ($w$ = widest level).
\end{complexitybox}

\solstep{Pitfalls and Variants}
\begin{warningbox}[Edges vs nodes]
Subtracting one here (to answer in edges) or adding one to an edge-height solution is the
classic wrong answer. Re-read whether the problem counts nodes or edges before submitting.
\end{warningbox}
Variant: \emph{Minimum Depth} (\probref{c01-min-depth}) below; \emph{Diameter of Binary
Tree} (LC 543) uses the same bottom-up idea on two paths at once (\chref{c09}).

\begin{insightbox}
Key takeaway: max-depth is \thmref{c01-height-recursion} with the base case set to the
convention the problem asks for. Master the single recursion, then change only the
null-return between $0$ (nodes) and $-1$ (edges).
\end{insightbox}

\solvedproblem{c01-min-depth}{Minimum Depth of Binary Tree}{LeetCode 111}
  {https://leetcode.com/problems/minimum-depth-of-binary-tree/}{LC Easy L2}
  {recursion, the one-child trap}

\begin{problemstatement}[Minimum Depth]
Return the number of nodes on the shortest path from the root to the nearest \emph{leaf}
(a node with no children). The path must end at a real leaf.
\end{problemstatement}

\begin{figure}[H]
\centering
\begin{tikzpicture}[diagram, scale=0.85, transform shape]
  % T1: root with only a right child
  \node[tnode] (a1) at (0,0) {1};
  \node[tnode] (b1) at (1,-1.2) {2};
  \draw[tedge] (a1) -- (b1);
  \node[annot] at (0.5,0.8) {$T_1$: answer $2$};
  % T2: single node
  \node[tnode] (a2) at (3.2,0) {1};
  \node[annot] at (3.2,0.8) {$T_2$: answer $1$};
  % T3: both children
  \node[tnode] (a3) at (6,0) {1};
  \node[tnode] (l3) at (5,-1.2) {2};
  \node[tnode] (r3) at (7,-1.2) {3};
  \draw[tedge] (a3) -- (l3); \draw[tedge] (a3) -- (r3);
  \node[annot] at (6,0.8) {$T_3$: answer $2$};
\end{tikzpicture}
\caption{The three test trees. $T_1$ is the trap: a naive $\min$ over children (with the
null child scoring $0$) answers $1$ -- but the lone node $1$ is \emph{not} a leaf, the
real answer is $2$. Every wrong submission dies on $T_1$.}\label{fig:c01-min-depth-trap}
\end{figure}

\solstep{Understanding the Problem}
The trap: a naive $\min(\text{left},\text{right})$ with $\depth(\nil)=0$ reports depth $1$
for a root that has only a right child, because the missing \emph{left} (a null, not a
leaf) looks shortest. We must only compare branches that actually reach a leaf.

\solstep{Key Observations}
\begin{observation}
A null child is \emph{not} a leaf and must never win the $\min$.
\end{observation}
\begin{observation}
If a node has exactly one child, its minimum depth is $1$ plus that child's minimum
depth -- the $\min$ has only one candidate.
\end{observation}

\solstep{Solution}
Three cases: null returns $0$; a node missing one side recurses on the side it has; a node
with both children takes $1+\min$ of the two.

\solstep{Proof of Correctness}
\begin{proofbox}[Induction on $\size(v)$]
For $v=\nil$ the answer is $0$. If $v$ has exactly one child $c$, the only root-to-leaf
paths go through $c$, so the minimum is $1+\depth_{\min}(c)$, which the code returns. If $v$
has two children, every leaf lies under one of them, and the shortest path picks the
smaller of the two correct recursive answers plus $v$. In the remaining case ($v$ is a
leaf) both children are null, the two-child formula gives $1+\min(0,0)$ only if we let
nulls score $0$ -- which is why we special-case missing children above instead.
\end{proofbox}

\solstep{Implementation}
\begin{lstlisting}[caption={Minimum Depth (LeetCode 111)}, label={lst:c01-min-depth-soln}]
#include <bits/stdc++.h>
using namespace std;

struct TreeNode {
    int val;
    TreeNode *left, *right;
    TreeNode(int x = 0) : val(x), left(nullptr), right(nullptr) {}
};

class Solution {
public:
    int minDepth(TreeNode* root) {
        if (root == nullptr) return 0;
        if (root->left  == nullptr) return 1 + minDepth(root->right);  // only right
        if (root->right == nullptr) return 1 + minDepth(root->left);   // only left
        return 1 + min(minDepth(root->left), minDepth(root->right));   // both
    }
};
\end{lstlisting}

\solstep{Complexity}
\begin{complexitybox}
Time $\bigO{n}$, memory $\bigO{h}$ recursion stack ($\bigO{n}$ worst case). A BFS that
returns as soon as it pops a leaf is often faster in practice (it stops at the first leaf)
and is naturally iterative.
\end{complexitybox}

\solstep{Pitfalls and Variants}
\begin{warningbox}[The one-child case is the whole problem]
Collapsing to $\min(\text{left},\text{right})$ with nulls scoring $0$ fails the moment a
node has a single child -- a huge fraction of the tests. Keep the two guard lines.
\end{warningbox}
Variant: minimum depth to \emph{any} node with a property (e.g. the nearest leaf matching
a predicate) is the same BFS-with-early-exit pattern.

% ---------------- stretch solved problem: two quantities, one DFS ----------------
\begin{attemptgate}
\textbf{LC 2265} (next) asks for something you have never computed: "nodes whose
value equals the average of their subtree." Try it -- the wall you hit is the
\emph{reason} bottom-up pairs exist. Give yourself 15 minutes.
\peek{You are already computing $\size(v)$ in one DFS. What is the second quantity
the average needs?}
\nudge{Return a \emph{pair} (sum, count) from each subtree, like
\coderef{c01-measure} returns height and size. Integer division truncates toward
zero -- does that ever mis-round a true average?}
\fullhint{Floor-division kills you: $7/2 = 3$ but the true mean is $3.5$. Check
$\mathrm{sum} = \mathrm{count} \times \mathrm{val}$ exactly -- the node counts iff its
value times the size equals the subtree sum.}
\end{attemptgate}

\solvedproblem{c01-count-nodes-avg}{Count Nodes Equal to Average of Subtree}{LeetCode 2265}
  {https://leetcode.com/problems/count-nodes-equal-average-of-subtree/}{LC Medium L3}
  {bottom-up pair return, one DFS, truncation trap}

\begin{problemstatement}[Count Nodes Equal to Average of Subtree]
Given the root of a binary tree, count nodes whose value equals the
\emph{floor of the average} of the values in their subtree (the node itself included).
$1 \le n \le 1000$, $1 \le \mathrm{val} \le 10^5$.
\end{problemstatement}

\solstep{Understanding the Problem}
For every node we need its whole subtree's sum and size, then compare
$\lfloor \mathrm{sum}/\mathrm{size} \rfloor$ with the node's value. "For every node, a
subtree statistic" is the signature of the bottom-up return -- exactly the machinery of
\coderef{c01-measure} with richer payloads.

\solstep{Key Observations}
\begin{observation}
$\depth$-style top-down passes cannot answer this: a node's subtree average is built
from its \emph{descendants}, so information must flow \emph{up}.
\end{observation}
\begin{observation}
Sum and size are both foldable in one DFS: $\mathrm{sum}(v) = \mathrm{val}(v) +
\sum_c \mathrm{sum}(c)$, $\size(v) = 1 + \sum_c \size(c)$ -- the pair returns together.
\end{observation}

\solstep{Think It Through With Me}
\begin{thinkbox}
\deadend{one DFS to fill arrays, then a second loop averaging each subtree by
re-walking it}{that is $O(n^2)$ on a chain -- passes for $n \le 1000$, but it teaches
nothing and dies at $10^5$.}
\deadend{computing $\mathrm{sum}$ bottom-up but doing the division at the end
($\mathrm{sum}/\mathrm{size}$ with ints)}{floor division: a subtree $\{3,4\}$ gives
$7/2 = 3$, and if the root value is $3$ you count it -- but $\lfloor\!\text{mean}\!\rfloor$
\emph{is} 3, so here it is right; the real killer is thinking $\mathrm{val} = \mathrm{sum}/\mathrm{size}$
is equivalent to $\mathrm{val} \cdot \mathrm{size} = \mathrm{sum}$. It is \emph{not}:
$\{3,4,4\}$ has mean $3.67$, floor $3$, yet $3 \cdot 3 \ne 11$.}
\aha{return $\mathrm{pair}(\mathrm{sum},\mathrm{size})$ up the recursion and test
$\mathrm{sum} / \mathrm{size} == \mathrm{val}$ using the \emph{same} integer division the
problem defines. The division is part of the \emph{specification}, not our shortcut --
so the floor question vanishes.}
\end{thinkbox}

\solstep{The Approach Ladder}
\begin{table}[H]\centering\small
\caption*{Approach ladder -- Count Nodes Equal to Average of Subtree}
\begin{tabular}{@{}c p{0.235\linewidth} p{0.28\linewidth} l c@{}}
\toprule
& \textbf{Approach} & \textbf{Why it wins / dies} & \textbf{Cost} & \textbf{Verdict}\\
\midrule
$\lv2$ & recompute each subtree's average by walking it
  & correct but quadratic on chains ($n\ge10^5$ dies); no reusable skill
  & $O(n^2)$ & $\colorbox{warningred!12}{\textcolor{warningred}{\scriptsize\bfseries TLE}}$\\
\addlinespace
$\lv3$ & one DFS returning $\mathrm{pair}(\mathrm{sum},\mathrm{size})$; test at each node
  & both stats fold bottom-up for free
  & $O(n)$ / $O(h)$ & $\ac$\\
\addlinespace
$\lv4$ & same DFS, but test $\mathrm{val}\cdot\mathrm{size}=\mathrm{sum}$ "to avoid division"
  & silently changes the spec (floor vs exact mean) -- WA on $\{3,4,4\}$
  & $O(n)$ & $\colorbox{warningred!12}{\textcolor{warningred}{\scriptsize\bfseries WA}}$\\
\bottomrule
\end{tabular}
\end{table}

\solstep{Solution}
One post-order DFS. Each call returns $\mathrm{pair}(\mathrm{sum},\mathrm{size})$ of its
subtree; on the way back up, after both children have returned, the node computes its own
pair, and increments the global answer when $\mathrm{sum}/\mathrm{size}$ (integer
division, exactly as the judge does it) equals its value.

\solstep{Proof of Correctness}
\begin{proofbox}[Induction on $\size(v)$]
For $v=\nil$ the returned pair is $(0,0)$, correct. For $v$ non-nil, by the induction
hypothesis the two children return the correct $(\mathrm{sum},\mathrm{size})$ of their
subtrees; the node adds its own value and one node, giving the correct pair for
$\subtree(v)$. The node is counted exactly when the spec's condition
$\lfloor\mathrm{sum}/\mathrm{size}\rfloor=\mathrm{val}$ holds, which the code tests with
the same integer division the judge uses. Summing these per-node decisions yields the
required count.
\end{proofbox}

\solstep{Implementation}
\begin{lstlisting}[caption={Count Nodes Equal to Average (LeetCode 2265)}, label={lst:c01-count-avg}]
#include <bits/stdc++.h>
using namespace std;

struct TreeNode {
    int val;
    TreeNode *left, *right;
    TreeNode(int x = 0) : val(x), left(nullptr), right(nullptr) {}
};

class Solution {
    int ans = 0;
    // returns {subtree_sum, subtree_size}
    pair<long long,int> dfs(TreeNode* v) {
        if (!v) return {0, 0};
        auto [sl, cl] = dfs(v->left);
        auto [sr, cr] = dfs(v->right);
        long long sum = sl + sr + v->val;
        int   cnt = cl + cr + 1;
        if (sum / cnt == v->val) ++ans;   // floor div == the spec
        return {sum, cnt};
    }
public:
    int averageOfSubtree(TreeNode* root) {
        ans = 0;
        dfs(root);
        return ans;
    }
};
\end{lstlisting}

\solstep{Trace}
\begin{tracebox}
Tree $\{3,4,4\}$ (root $3$, two leaf children $4,4$): each leaf returns $(4,1)$ and is
counted ($4/1=4$). The root sums to $3+4+4=11$, count $3$, $11/3=3$ (integer) $=\;$
root value, so it is counted too -- answer $3$. Note the WRONG
$\mathrm{val}\cdot\mathrm{size}=\mathrm{sum}$ test would reject the root ($3\cdot3\ne11$).
\end{tracebox}

\solstep{Complexity}
\begin{complexitybox}
Time $\bigO{n}$ -- one visit per node, $O(1)$ combine. Memory $\bigO{h}$ stack. Sum in
\code{long long}: at $n=1000$, $\mathrm{val}\le 10^5$ the bound is $10^8$ -- safe even
in \code{int}, but the habit is what saves you at $n=10^5$ in the next problem.
\end{complexitybox}

\solstep{Verification and Edge Cases}
\begin{itemize}[nosep]
  \item single node $\{5\}$: sum/size $=5$ -- counted. \hfill\lv{3}
  \item chain of equal values: every subtree averages to that value -- answer $n$.
  \item floor trap from the thinkbox: $\{3,4,4\}$ root 3, mean $11/3$ floor $3$ -- counted
        (integer div is the \emph{spec}); the WRONG $\mathrm{val}\cdot\mathrm{size}=\mathrm{sum}$
        test rejects it -- that is the $\lv4$ WA rung.
\end{itemize}

\solstep{Pitfalls and Variants}
\begin{warningbox}[Truncation is the specification]
Do not "improve" the floor division into an exact-mean test. The problem defines the
average as integer floor; matching the judge's rounding is part of being correct.
\end{warningbox}

\begin{insightbox}
Key takeaways: (1) "for every node, a subtree statistic" $\Rightarrow$ return a
\emph{pair} bottom-up in a single DFS; (2) the naive re-walk is $O(n^2)$ -- the pair
return makes it $O(n)$; (3) reproduce the spec's rounding, do not substitute a "nicer"
test. This pair-return is the doorway to tree DP (\chref{c10}).
\end{insightbox}

\begin{hardeningbox}
\textbf{$\lv4$ stress:} $n$ up to $10^5$ as a chain -- the $\bigO{h}$ stack overflows on
many judges. Convert to the explicit-stack post-order twin (pairs stored in a parallel
array, exactly the \coderef{c01-measure} bookkeeping). \textbf{$\lv5$ variant:} same pair
return but the average is over the \emph{path to root} (top-down accumulator, no pairs
needed) -- practise flipping direction on purpose. \textbf{$\lv5$ stretch:} count nodes
equal to the \emph{median} of their subtree -- pairs are no longer enough; that is
DSU-on-tree territory (\chref{c08}).
\end{hardeningbox}

\subsection{Practice Problems}

\begin{stuckbox}
\begin{itemize}[nosep, leftmargin=1.4em]
  \item Test your base case on a \emph{single node}: is the answer $0$ (edges) or $1$
        (nodes)? Convert conventions explicitly (\ref{sec:c01-anatomy} convention note).
  \item ``Minimum depth'' on a one-child node: a null child is \emph{not} a leaf --
        ignoring this gives the classic wrong answer.
  \item Skewed tree of $10^5$ nodes: will your recursion survive? If not, convert to an
        explicit stack before debugging anything else.
  \item For diameter-style problems, separate the two quantities: what you \emph{return}
        (best downward path) vs what you \emph{update globally} (best path through a node).
\end{itemize}
\end{stuckbox}

\begin{practiceset}
  \practice{Sum of Left Leaves}{LeetCode 404}
    {https://leetcode.com/problems/sum-of-left-leaves/}{LC Easy L2}
    {recurse, passing a flag for "I am a left child"; add the value when a leaf is left.}
  \practice{Tree Diameter (Longest Path)}{CSES 1131}
    {https://cses.fi/problemset/task/1131}{CSES L3}
    {two BFS/DFS sweeps: farthest node from an arbitrary start, then farthest from that. See \chref{c09}.}
  \practice{Diameter of Binary Tree}{LeetCode 543}
    {https://leetcode.com/problems/diameter-of-binary-tree/}{LC Easy L2}
    {bottom-up: return height, update a global best with left-height + right-height.}
  \practice{Binary Tree Longest Consecutive Sequence}{LeetCode 298}
    {https://leetcode.com/problems/binary-tree-longest-consecutive-sequence/}{LC Medium L3}
    {Stretch: carry the current run length down in a DFS; compare both children, do not branch-max naively.}
\end{practiceset}

% ---------------- spaced-repetition card ----------------
\begin{recallbox}
\selftest{
  \item With the edges convention, $\height(\nil)=$ \underline{\hspace{0.9cm}} and a lone
        node has height \underline{\hspace{0.9cm}}.
  \item Which quantities flow \emph{down} the recursion (carried in) and which bubble
        \emph{up} (returned)?
  \item Min-depth dies on which specific tree shape? Draw it.
  \item In LC 2265 why is testing $\mathrm{val}\cdot\mathrm{size}=\mathrm{sum}$ a WA
        against the floor-average spec?
  \item Write the one-DFS $\langle\height,\size\rangle$ measurement skeleton from memory,
        base case first.
}
One DFS fills depth, height and size by pushing depth in and returning
$\langle\height,\size\rangle$ out.
\end{recallbox}
